% Standalone problem document, generated from the adjacent README.md.
% Compile with XeLaTeX. Mathematical content is maintained in Markdown.
\documentclass[11pt,a4paper]{article}
\usepackage[fontsize=10.5pt]{scrextend}
\usepackage[margin=22mm,headheight=15pt,headsep=9mm,footskip=11mm]{geometry}
\usepackage{fontspec}
\setmainfont{texgyrepagella}[Extension=.otf,UprightFont=*-regular,BoldFont=*-bold,ItalicFont=*-italic,BoldItalicFont=*-bolditalic]
\setsansfont{texgyreheros}[Extension=.otf,UprightFont=*-regular,BoldFont=*-bold,ItalicFont=*-italic,BoldItalicFont=*-bolditalic]
\setmonofont{texgyrecursor}[Extension=.otf,UprightFont=*-regular,BoldFont=*-bold,ItalicFont=*-italic,BoldItalicFont=*-bolditalic,Scale=MatchLowercase]
\usepackage{amsmath,amssymb}
\usepackage{unicode-math}
\setmathfont{texgyrepagella-math.otf}
\usepackage{xcolor,microtype,enumitem,needspace,fancyhdr}
\usepackage{bookmark}
\definecolor{ink}{HTML}{17344B}
\definecolor{accent}{HTML}{197D80}
\definecolor{muted}{HTML}{596773}
\hypersetup{colorlinks=true,linkcolor=accent,urlcolor=accent,pdftitle={MF-02 - Multiplication
overhead of cubic sign
compositions},pdfauthor={Open Problems in Numerical Linear Algebra}}
\urlstyle{same}
\setlength{\parindent}{0pt}
\setlength{\parskip}{5pt plus 1pt minus 1pt}
\linespread{1.04}
\setlength{\emergencystretch}{3em}
\widowpenalty=10000
\clubpenalty=10000
\displaywidowpenalty=10000
\allowdisplaybreaks[1]
\setlist{nosep,leftmargin=1.5em}
\providecommand{\tightlist}{\setlength{\itemsep}{0pt}\setlength{\parskip}{0pt}}
\providecommand{\pandocbounded}[1]{#1}
\pagestyle{fancy}
\fancyhf{}
\fancyhead[L]{\sffamily\footnotesize\color{muted}OPEN PROBLEMS IN NUMERICAL LINEAR ALGEBRA}
\fancyhead[R]{\sffamily\bfseries\footnotesize\color{accent}MF-02}
\fancyfoot[L]{\parbox[b]{0.9\textwidth}{\sffamily\fontsize{8}{10}\selectfont\color{muted}Editorial ratings; a bounded search cannot certify that a problem remains unsolved.\\Literature check: 2026-09-22}}
\fancyfoot[R]{\sffamily\footnotesize\color{muted}\thepage}
\renewcommand{\headrulewidth}{0.3pt}
\renewcommand{\headrule}{\hbox to\headwidth{\color{accent}\leaders\hrule height \headrulewidth\hfill}}
\makeatletter
\renewcommand{\section}{\Needspace{5\baselineskip}\@startsection{section}{1}{0pt}{12pt plus 2pt minus 2pt}{4pt}{\sffamily\bfseries\large\color{ink}}}
\renewcommand{\subsection}{\Needspace{4\baselineskip}\@startsection{subsection}{2}{0pt}{9pt plus 1pt minus 1pt}{3pt}{\sffamily\bfseries\normalsize\color{ink}}}
\makeatother
\setcounter{secnumdepth}{0}
\begin{document}
{\sffamily\small\color{muted}Matrix functions and stability\par}
\vspace{3pt}
{\raggedright\hyphenpenalty=10000\exhyphenpenalty=10000\sffamily\bfseries\fontsize{21}{24}\selectfont\color{ink}Multiplication
overhead of cubic sign compositions\par}
\vspace{7pt}
{\sffamily\small\textbf{Difficulty:} challenging\quad\textbf{Importance:} interesting
to the community\par}
{\sffamily\small\bfseries\color{accent}Status: Lean verified\par}
\vspace{4pt}
{\color{accent}\hrule height 0.8pt}
\vspace{6pt}
\textbf{Rating rationale:} Challenging because optimal cubic
compositions must be compared with unrestricted evaluation programs;
community impact comes from reusable matrix-sign iterations.

\subsection{Lean verification --- 22 September
2026}\label{lean-verification-22-september-2026}

The full canonical uniform stage bound is formally proved for every
multiplication budget and every real gap in (0,1), including both
small-budget endpoints. Seven exports establish actual register-program
costs, genuine approximation errors and least feasible stage counts.
They passed LeanCert kernel audits, two independent nonauthor final
reviews, and
\href{https://github.com/sidneyholden1/OpenProblemsInNLA/actions/runs/35693827256}{fresh
isolated Linux Comparator/default-kernel verification} with rejection
and sandbox controls.

\href{https://github.com/sidneyholden1/OpenProblemsInNLA/blob/main/matrix-functions-and-stability/MF-02/lean/README.md}{Proof
and reviews} ·
\href{https://github.com/sidneyholden1/OpenProblemsInNLA/tree/9aba1c7fed69ad95023b217a28cb486a7b9983ef/matrix-functions-and-stability/MF-02/lean}{Immutable
proof revision} ·
\href{https://github.com/sidneyholden1/OpenProblemsInNLA/blob/main/docs/lean/verification-2026-09-22/MF-02/README.md}{Retained
Linux evidence}. Exposition: George Stepaniants; prior asymptotic order:
Cheon--Kim--Kim; cubic construction: Chen--Chow. Formalization: Sidney
Holden with OpenAI Codex assistance. This verifies the stated uniform
asymptotic order; no novelty, exact optimum, same-budget optimality or
external human peer review is claimed.

\subsection{Resolution - uniform constant-factor asymptotic
order}\label{resolution---uniform-constant-factor-asymptotic-order}

\textbf{Resolution recorded 2026-09-12.} The canonical asymptotic-order
target is settled:

\[T_{\min}(m,\delta)=\Theta(m+1)\qquad(0<\delta<1),\]

with absolute constants uniform in the gap, including gaps depending on
the multiplication budget. The
\href{https://github.com/ajt60gaibb/OpenProblemsInNLA/blob/main/matrix-functions-and-stability/MF-02/solution.md}{Theorem
in Section 1 and proof in Sections 2-5} give the explicit bounds

\[\left\lfloor\frac m2\right\rfloor\le T_{\min}(m,\delta)\le m\quad(m\ge2),
\qquad T_{\min}(0,\delta)=T_{\min}(1,\delta)=1.\]

\textbf{Expository proof-note author:} George Stepaniants, Department of
Computing and Mathematical Sciences, California Institute of Technology,
Pasadena, California, USA.
\href{https://github.com/ajt60gaibb/OpenProblemsInNLA/blob/main/matrix-functions-and-stability/MF-02/solution.pdf}{Proof
PDF} ·
\href{https://github.com/ajt60gaibb/OpenProblemsInNLA/blob/main/matrix-functions-and-stability/MF-02/solution.tex}{Standalone
TeX}.

\textbf{Prior work and exact scope.} Cheon, Kim and Kim's
\href{https://doi.org/10.1007/978-3-030-64834-3_8}{ASIACRYPT 2020 work},
especially Lemma 3 and its constant-factor complexity analysis, already
yields uniform constant-factor order by combination with the classical
degree bound. The optimized cubic is the earlier
\href{https://www.mcs.anl.gov/papers/P5059-0114.pdf}{Chen-Chow
construction, equations (3.3)-(3.6)}, also discussed in
\href{https://arxiv.org/html/2505.16932v5}{Polar Express, Appendix F}.
The present self-contained note supplies the explicit comparison
\(T_{\min}\le m\); it claims neither a new cubic iteration nor first
discovery of constant-factor optimality. The exact smallest stage count,
the optimal leading constant, and the stronger source question comparing
the errors at the same multiplication budget remain unanswered.

\textbf{Verification.} A separate
\href{https://github.com/ajt60gaibb/OpenProblemsInNLA/blob/main/references/stepaniants-mf02-2026-09-12/verification/independent-review/MF-02-independent-review.md}{independent
Codex-agent mathematical and scope review} returned PASS for the theorem
and the literal canonical asymptotic-order target. A second
\href{https://github.com/ajt60gaibb/OpenProblemsInNLA/blob/main/references/stepaniants-mf02-2026-09-12/verification/source-review/REVIEW.md}{independent
source and attribution review} confirms the scope and prior-work
distinctions. The
\href{https://github.com/ajt60gaibb/OpenProblemsInNLA/blob/main/references/stepaniants-mf02-2026-09-12/README.md}{submission
record} preserves both reviews, the original candidate and exact checks.
Substantial AI assistance is disclosed. This is informal automated-agent
review, not external human peer review or formal verification. No
novelty or priority certification is asserted.

The original target, permanent ID, references and dated history remain
below. The displayed ratings are retained as historical ratings of that
target.

\newpage

\subsection{Context and notation}\label{context-and-notation}

For \(m\in\mathbb N_0\) and \(0<\delta<1\), put

\[I_\delta=[-1,-\delta]\cup[\delta,1].\]

Let \(\mathcal P_m\) consist of real polynomials computed from \(1,x\)
by straight-line programs using at most \(m\) nonscalar multiplications;
real linear combinations cost nothing. Define

\[E_m(\delta)=\inf_{p\in\mathcal P_m}
\max_{x\in I_\delta}|p(x)-\mathop{\mathrm{sign}}\nolimits(x)|.\]

The arithmetic model concerns a single polynomial identity valid for
matrices of every size.

\subsection{Problem statement}\label{problem-statement}

Define

\[C_T(\delta)=\inf_{a_t,b_t\in\mathbb R}
\max_{x\in I_\delta}|(q_T\circ\cdots\circ q_1)(x)-\mathop{\mathrm{sign}}\nolimits(x)|,
\qquad q_t(x)=a_tx+b_tx^3.\]

Determine the asymptotic dependence on \(m,\delta\) of

\[T_{\min}(m,\delta)=\inf\{T\in\mathbb N_0:C_T(\delta)\le E_m(\delta)\},\]

where the empty composition is \(x\) and \(\inf\varnothing=+\infty\).
Each stage costs at most two matrix products.

\subsection{References and status
evidence}\label{references-and-status-evidence}

Amsel et al., \href{https://arxiv.org/html/2602.05394v3}{Simons workshop
report}, §6.3, Problem 6.5. Rubensson, Jarlebring and Lorentzon,
\href{https://arxiv.org/html/2606.24701v1}{degree-eight recursive
expansion}, §7, provides a June 2026 follow-up discussion; it does not
settle this comparison.

\subsection{Scope}\label{scope}

\href{https://github.com/ajt60gaibb/OpenProblemsInNLA/blob/main/matrix-functions-and-stability/MF-01/README.md}{MF-01}
asks for an optimum value over general evaluation programs. This problem
measures the price of a particular, reusable composition architecture;
the two tasks are related but have different requested outputs.

\subsection{Audit --- 2026-09-10}\label{audit-2026-09-10}

Rechecked \href{https://arxiv.org/html/2602.05394v3\#S6.SS3}{Problem
6.5} and the \href{https://arxiv.org/html/2606.24701v1}{degree-eight
paper's concluding discussion}. The optimal composition comparison
remains posed. Searches for cubic and recursive sign-expansion
improvements found no resolution of this exact multiplication-overhead
target.
\end{document}
